\section{Related Work}
\label{sec:related-work}

Our method sits between schedule optimization and solver design. The closest
prior work asks how to choose the discretization grid; our distinction is that
we estimate the local oracle-start defect of the target solver under the
teacher marginal and then convert that defect into an equal-mass monitor.

\textbf{Diffusion samplers and hand-designed schedules.}
DDPMs, DDIMs, and score-SDE samplers define reverse generative dynamics and
expose a finite set of model-evaluation locations
\citep{ho2020ddpm,song2022ddim,song2021sde}. EDM makes the design space of
preconditioning, noise parameterization, and polynomial noise schedules explicit
\citep{karras2022edm}. These schedules are cheap, stable, and reusable, which
makes them strong practical baselines. They are not, however, diagnostics of a
particular solver's local prediction defect.

\textbf{Offline and path-geometric schedule optimization.}
AYS is the strongest direct schedule-optimization baseline in our setting. It
uses stochastic calculus to build a Girsanov-based KL upper bound between the
true generative SDE and a solver-induced linearized SDE, then searches for a
schedule specific to a model, dataset, and solver \citep{sabour2024ays}. Its
published Stable Diffusion schedules are therefore important baselines in the
text-to-image sweep.

Score-Optimal Diffusion Schedules are especially close conceptually
\citep{williams2024scoreoptimal}. Williams et al. define a diffusion path cost
from the work a Langevin corrector must perform after a predictor step. The
resulting optimal generator is an inverse path-length parameterization: local
costs are equalized along the sampling path. We keep this equalization
principle but change the object being equalized. Instead of a path-intrinsic
score-corrector cost under perfect-score assumptions, we measure the
oracle-start one-step residual of the actual deterministic solver against a
teacher flow, aggregate it over $p_u$, and use the $q$-th root dictated by the
measured local defect order.

Other optimized-timestep methods also move beyond fixed grids. Xue et al.
optimize time steps by minimizing the distance between a ground-truth ODE
solution and the numerical solution induced by a specific solver
\citep{xue2024optimizedtimesteps}. This is close to the teacher-student
motivation of our method, but their method solves an offline constrained
optimization problem, whereas ours estimates a distributional local defect
profile and materializes the schedule by inverse-CDF equalization.

\textbf{Solver acceleration.}
PNDM, DPM-Solver, DPM-Solver++, UniPC, and related predictor-corrector solvers
reduce sampling cost by changing the numerical update rule
\citep{liu2022pndm,lu2022dpmsolver,lu2025dpmsolverpp,zhao2023unipc}. STORK
uses stabilized Taylor orthogonal Runge--Kutta ideas to improve mid-NFE
diffusion and flow-matching sampling while addressing stiffness and
structure-dependence \citep{tan2025stork}. These methods are complementary to
our method: they change how a step is computed, while ours changes where a
fixed solver evaluates the model. The distinction also matters for theory. Our
main guarantee is stated for deterministic oracle-start ODE defects; multistep,
stochastic, and scheduler-history solvers require either augmented-state
analysis or empirical framing.

\textbf{Learned fast samplers.}
Differentiable sampler search and progressive distillation reduce NFE by
optimizing sampler parameters or distilling many-step samplers into fewer-step
generators \citep{watson2022fastsamplers,salimans2022progressive}. Our method
does not update model weights and does not optimize an image-quality objective.
Negative cases therefore diagnose the schedule monitor and its calibration
distribution rather than a learned generator.
